\documentstyle[% 


righttag% 


,11pt% 


,amssymb% 


,russian%               remove in Eng. 


,izvru%                 Set izven% in Eng. 


% ,izvru-ks             for brief communications 


]{amsart} 


 


%       Math definitions 


 


\newcommand{\fs}{f\sup} 


 


\theoremstyle{plain} 


\theorembodyfont{\it} 


\newtheorem{propnonum}{\hskip\parindent Предложение} 


\renewcommand{\thepropnonum}{} 


 


\theoremstyle{definition} 


\theorembodyfont{\rm} 


\newtheorem{cornonum}{\hskip\parindent Следствие} 


\renewcommand{\thecornonum}{} 


 


%\hoffset=-15mm                  For Eng. 


 


\setcounter{page}{20} 


\god{1997} 


\nomer{8~(423)} %                Remove in Eng. 


%\nomer{Vol.\,41, No.\,8}        Set in Eng. 


% \str{75-81}                    Set in Eng. 


\udk{519.216} 


\author{\it А.И.\,Володин} 


\title{Сходимость взвешенных сумм случайных элементов\\ 


в банаховых пространствах типа $p$} 


 


\date{} 


% \pagestyle{myheadings}         Set in Eng. 


 


\pagestyle{plain} %              Remove in Eng. 


\begin{document} 


 


% \thispagestyle{plain}          Set in Eng. 


 


\maketitle 


 


% Body text 





\section{Введение}





Вводится понятие $\fs$-сходимости случайных


элементов, которая сильнее, чем сходимость п.\,н., и для этой


сходимости изучаются законы больших чисел для взвешенных сумм


случайных элементов со значениями в банаховом пространстве


типа $p$.





Более сильная, чем п.\,н. сходимость определяется следующим образом:


для возрастающей непрерывной функции


$f:\bold R^+ \to\bold R^+$,


$f(0)=0$, $f(\infty )=\infty$, и случайных элементов


$(T_n)^{\infty }_1$


со значениями в


банаховом пространстве $E$ последовательность $T_n\to 0$ в


смысле


$\fs$, если $\bold E f(\sup\limits_{k\ge n}\| T_k \|)\to 0$


при $n\to \infty$. Очевидно, $\fs$-сходимость


сильнее как сходимости почти наверное, так и сходимости


в пространстве Орлича $L_{f} (E)$ (если $f$ удовлетворяет


$\Delta_{2}$-условию (\cite{Kra}, теорема 9.4)).





В работе находятся моментные условия для норм случайных


элементов со значениями в пространстве типа $p$, при которых взвешенные


суммы этих элементов сходятся к нулю в смысле $\fs$.





Введем ряд определений и обозначений: $E$ --- вещественное


сепарабельное банахово пространство, $(X_{k})^{\infty }_{1}$ ---


последовательность независимых симметричных случайных элементов со


значениями в $E$, через $C$ будем обозначать любую положительную


постоянную, которая может принимать различные значения даже


в пределах одной формулы. Скажем, что последовательность


случайных элементов $(X_{k})$ {\it стохастически доминируется}


положительной


случайной величиной $\xi$ (обозначение $(X_{k} )\prec \xi )$, если


существует такое $C> 0$, что для всех $t>0$


$$


\sup_{k\ge 1} \bold P \{\| X_{k} \| >t\} \le C \bold P \{\xi >t\}.


$$


Используемые в дальнейшем понятия и свойства


радемахеровского и устойчивого типов банахова пространства можно найти в


(\cite{Pis}, ch.\,3--4). Напомним также


\medskip





{\bf Неравенство Квапеня} \cite{Sz}.


{\it Пусть $(X_k)_1^n$ --- независимые симметричные случайные 


элементы и $(s_k)^n_1\subset\bold R$. Тогда для всех $t\ge 0$ вероятность}


$$


\bold P\Bigl\{\Bigl\|\sum^n_{k=1}s_kX_k\Bigr\|>t\Bigr\}\le


2\bold P\Bigr\{\,\max_{1\le k\le n}|s_k|\Bigl\|\sum^n_{k=1}X_k\Bigr\|>t\Bigr\}.


$$ 


\medskip





Несколько слов об истории вопроса. Исследование сходимости взвешенных сумм


является естественным продолжением изучения закона больших чисел Марцинкевича.


Закон больших чисел Марцинкевича изучался T.A.\,Азларовым и Н.A.\,Володиным


\cite{Azl}, A.~дe~Акостой \cite{Aco} в банаховом пространстве 


радемахеровского


типа; M.\,Maркусом и В.A.\,Вoйчинским \cite{Marc} в банаховом пространстве


устойчивого типа. Обобщения этих результатов на взвешенные суммы случайных


элементов были получены А.\,Адлером, А.\,Росальским и Р.Л.\,Тейлором


\cite{Adl},


Т.\,Микошем и Р.\,Норвайшей \cite{Mik}, Р.Л.\,Тейлором \cite{Tay}.


Данная работа продолжает исследования по этой тематике.


Ряд результатов этой статьи анонсировался в \cite{Vol1} и \cite{Vol2}.





\section{Формулировка основного результата и доказательство лемм}





Пусть $(X_{k})$ --- последовательность независимых симметричных случайных 


элементов со


значениями в банаховом пространстве $Е$,


${\bold a}=\{a_{k} (n),\;n\in \bold N,\;1\le k\le n\}$ ---


двойная последовательность вещественных чисел, которая в дальнейшем будет


называться {\it взвесью}


(естественно, не все $a_{k} (n)=0)$.


Изучается


$\fs$-сходимость к


нулю взвешенных сумм вида


$$


T_{n} = \sum^{n}_{k=1} a_{k} (n) X_{k}.


$$.%%%%%%%<-- remove dot in Eng.





Обозначим $\varphi (n)=1/\max\limits_{1\le k\le n}|a_{k}(n)|$,


$n\in\bold N$, и потребуем, чтобы


последовательность $\varphi =(\varphi (n))$ возрастала и


$\lim\limits_{n\to \infty}\varphi (n)= \infty $.


Так как мы изучаем закон больших чисел в пространствах


типа $p$, то необходимы некоторые условия роста как функции $f$


(из определения $\fs$-сходимости), так и последовательности


$\varphi$


(coответствующей взвеси ${\bold a}$), cвязанные с величиной $p$.


Введем следующие классы непрерывных функций:


$$


{\bold I}_{p} =\{f:\bold R^{+} \to\bold R^{+},\quad


f(0)=0,\quad f(\infty )=\infty,\quad f\text { возрастает и }


\int^{\infty }_{1} f(t)t^{-p-1} dt< \infty \}


$$


и ${\bold G}_{p}$ --- класс таких полуаддитивных функций $f$,


что существует


$q<p$ для которого $f\in {\bold I}_{q} $. Напомним, что полуаддитивность


функции $f$ означает $f(t+s)\le C(f(t)+f(s))$ для некоторой


постоянной $C$, зависящей только от $f$, и всех $t>0$, $s>0$.





Скажем, что взвесь ${\bold a}\in F_{p} $, если соответствующая


последовательность $\varphi $ является правильно меняющейся


последовательностью с параметром


$1/p$, $1<p<2$, т.\,е. $\varphi (n)= n^{1/p} L(n)$, где


$\lim\limits_{n\to\infty} L(mn)/ L(n)= 1$ для всех


$m\in \bold N$ (\cite{Sen}).





Основной результат статьи составляет





\begin{thm} Если банахово пространство $Е$ имеет устойчивый тип


$p$, $1<p<2$, взвесь ${\bold a}\in F_{p} $, функция $f\in {\bold G}_{p}$,


$(X_{k})$ --- независимые симметричные случайные элементы и


$(X_{k} )\prec \xi $, то из


$\sum\limits^{\infty }_{k=1} \bold P \{\xi >\varphi (k)\}<\infty$


следует $T_{n} \to 0$ в $\fs$ при $n\to \infty$.


\end{thm}





Для доказательства теоремы нам понадобится ряд лемм.


Первая лемма является простым обобщением известного


классического результата (см. \cite{Sto}, c.\,127--128;


\cite{Choi}, лемма 2.2),


поэтому приводится без доказательства.





\begin{lem} Пусть $(X_{k} )\prec \xi $ и


$\sum\limits^{\infty }_{k=1} \bold P \{\xi >\varphi (k)\}< \infty $.





{\em 1)} Если для некоторого


$r>0\;$


$\sum\limits^{\infty }_{k=n} \varphi ^{-r} (k)= O(n\varphi ^{-r} (n))$,


то


$$


\sum^{\infty }_{n=1} \varphi ^{-r} (n)\bold E


\| X_{n} \|^{r} I\{\| X_{n} \| <\varphi (n)\}< \infty.


$$





{\em 2)} Если для некоторого


$q>0\;$


$\sum\limits^{n}_{k=1} \varphi ^{-q} (k)= O(n\varphi ^{-q} (n))$, и


существует такое $C>0$, что $\varphi (k+1)\le C\varphi (k)$ для всех


$k\in \bold N$, то


$$


\sum^{\infty }_{n=1} \varphi ^{-q} (n)\bold E


\| X_{n} \|^{q} I\{\| X_{n} \|


\ge \varphi (n)\}< \infty.


$$


\end{lem}





\begin{lem} Если $(X_{k} )$ --- независимые симметричные случайные


элементы, то


$\bold E f(\sup\limits_{n\ge k} \| T_{n} \|)\le


4\bold E f(\sup\limits_{n\ge k} \| S_{n} \| /\varphi (n))$


для любого $k\ge 1;$ $S_n=\sum\limits^n_{k=1}X_k$.


\end{lem}





{\bf Доказательство.} По неравенству Леви для всех $k\le m$


$$


\bold Q=\bold P \{\sup_{k\le n\le m}\| T_{n}\| >


\varepsilon \}\le 2\bold P \{\| T_{m} \| >\varepsilon \}.


$$


Тогда по неравенству Квапеня


$$


\bold Q \le 4{\bold P}


\{\| S_{m} \| /\varphi (m)>\varepsilon \}


\le 4{\bold P} \{ \sup_{k\le n\le m} \| S_{n} \| /


\varphi (n)>\varepsilon \}.


$$


Следовательно,


${\bold P} \{\sup\limits_{n\ge k}\| T_{n} \|>


\varepsilon \}\le 4{\bold P} \{\sup\limits_{n\ge k}\| S_{n}\|


/\varphi (n)>\varepsilon \}$. Пусть $g(t)$ ---


обратная к $f(t)$ функция: $f(g(t))=t$.


Тогда


\begin{multline*}


{\bold E} f(\sup_{n\ge k}\|T_{n} \|)


=\int_0^{\infty}{\bold P} \{f(\sup_{n\ge 1}\|T_{n} \|)>t\}dt= 


\int_0^{\infty}{\bold P} \{\sup_{n\ge 1}\|T_{n} \|>g(t)\}dt\le\\


%


\le 4\int_0^{\infty}{\bold P} \{\sup_{n\ge 1}\|S_{n} \|/\varphi (n)>g(t)\}dt


=4{\bold E} f(\sup_{n\ge k}\|S_{n} \|/\varphi(n)).\quad\square


\end{multline*}





Третья лемма есть аналог результата В.Д.\,Чои и С.Х.\,Санга


(\cite{Choi}, теорема 2.1). Существенно новым является рассмотрение


случая банахова пространства и нормировки общего вида.





\begin{lem} Пусть $E$ имеет радемахеровский тип $p$, $1<p<2$.


Если $(X_{k})$ --- такие независимые симметричные случайные


элементы, что


$$


\sum^{\infty }_{k=1} {\bold E} \|X_{k}\|^{p} /\varphi ^{p} (k)<\infty,


$$


то для любой функции $f\in {\bold I}_{p} $ среднее


${\bold E} f(\sup\limits_{n\ge 1}\|S_{n} \|/\varphi (n))<\infty $.


\end{lem}





{\bf Доказательство.} Пусть $g(t)$ --- обратная к $f$ функция:


$f(g(t))= t$, $t\ge 0$. Для любого натурального $k\ge 0$ положим


$m_{k} =\min \{n:\varphi (n)\ge 2^{k} \}$ и


$N_{k} =\{n: m_{k} \le n<m_{k+1} \}$. Тогда в силу неравенства


Леви


\begin{multline*}


{\bold Q} = {\bold P} \{f(\sup_{n\ge 1}\|S_{n} \|


/\varphi (n))>t\}={\bold P} \{\sup_{k\ge 0}


\max_{n\in N_k} \|S_{n} \|/\varphi (n)>g(t)\}\le \\


%


\le \sum^{\infty }_{k=0}{\bold P} \{\max_{n\in N_k}


\|S_{n} \|/\varphi (n)> g(t)\}


\le \sum^{\infty }_{k=0} {\bold P} \{\max_{n\in N_k}


\|S_{n} \|> 2^{k} g(t)\}


\le 2\sum^{\infty }_{k=0} {\bold P} \{\|S_{m_{k+1} -1}\|> 2^{k} g(t)\}.


\end{multline*}


Теперь применим неравенство Чебышева и воспользуемся тем,


что пространство $E$ имеет тип $p$ (\cite{Pis}, ch.\,3):


$$


{\bold Q} \le 8g^{-p} (t)\sum^{\infty }_{k=0} 2^{-kp}


{\bold E} \|S_{m_{k+1} -1} \|^{p} \le Cg^{-p} (t)


\sum^{\infty }_{k=0} 2^{-kp}\sum^{m_{k+1} -1}_{i=1}


{\bold E}\|X_{i} \|^{p}.


$$


Поменяв порядок суммирования, получим


$$


{\bold Q} \le Cg^{-p} (t)\sum^{\infty }_{i=1}


{\bold E} \| X_{i} \| \sum_{k\ge k_{i} }


2^{-kp},


$$


где $k_{i} =\min \{k: m_{k+1} -1\ge i\}$. Для этого


$k_{i} $ справедливы неравенства


$m_{k_{i} +1} -1\ge i$ и $\varphi (m_{k_{i} +1} -1)> 2^{k_{i} +1} $.





Оценим сумму


$$


\sum_{k\ge k_{i} } 2^{-kp} =\frac{2^p}{1 - 2^{-p} }


2^{-(k_{i} +1)p}


\le C(\varphi (m_{k_{i} +1} -1))^{-p} \le C\varphi ^{-p} (i).


$$





Итак, ${\bold Q} \le Cg^{-p} (t)\sum\limits^{\infty }_{i=1}


{\bold E} \| X_{i} \| ^{p} $/$\varphi ^{p} (i)$.





Сделав замену переменных


$t=f(s)$ в интеграле $\int\limits_1^{\infty}f(s)ds^{-p} $,


который сходится по условию, и проинтегрировав по частям, получим


$\int\limits_1^{\infty}g^{-p} (t)dt<\infty $. Но тогда


$$


{\bold E}


f(\sup_{n\ge 1}\| S_{n} \| /\varphi (n))


\le 1 +\int_1^{\infty}


{\bold P} \{f(\sup_{n\ge 1}


\| S_{n} \| /\varphi (n)>t\}dt 


\le 1 + C\sum^{\infty }_{i=1}


{\bold E} \| X_{i} \| ^{p}/


\varphi ^{p} \int_1^{\infty} g^{-p} (t)dt<\infty.\quad\square


$$





\section{Основные результаты}





Следующее предложение может рассматриваться как закон больших чисел


Чжуна (\cite{Hof}, тeoрeма 3.1) относительно $\fs$-сходимости


для взвешенных сумм.





\begin{propnonum} Пусть заданы взвесь ${\bold a}$ и соответствующая


последовательность $\varphi $ и пространство $E$ имеет радемахеровский


тип $p$, $1<p<2$. Если $(X_{k} )$ --- такие независимые симметричные


случайные элементы, что $\sum\limits^{\infty }_{k=1}


{\bold E} \| X_{k} \| ^{p} /\varphi ^{p} (k)<\infty$,


то для любого


$f\in {\bold I}_{p} $ последовательность $T_{n} \to 0$ в $\fs$.


\end{propnonum}





\begin{pf} Так как пространство $E$ имеет тип $p$ (\cite{Pis},


ch.\,3), то


$$


{\bold E}


\Bigl\|\, \sum^{\infty }_{k=1} X_{k} /\varphi (k)\Bigr\| ^{p}


\le C\sum^{\infty }_{k=1}


{\bold E} \| X_{k} \| ^{p} /\varphi ^{p} (k),


$$


т.\,е. ряд


$\sum\limits^{\infty }_{k=1} X_{k} /\varphi (k)$


сходится в $L^{p}(Е)$. Но тогда он сходится п.\,н. (\cite{Tay},


\S\,4.5). По лемме Кронекера


$\frac{1}{\varphi (n)}\sum\limits^n_{k=1}X_k\to 0$


п.\,н. при


$n\to\infty $.





Применив неравенство Квапеня, находим


$$


\bold P\{\|T_m\|>t\}\le 2\bold P\Bigl\{\frac{1}{\varphi(n)}


\Bigl\|\sum^n_{k=1}X_k\Bigr\|>t\Bigr\}\to 0.


$$


т.\,е. $T_n\to 0$ по вероятности при $n\to\infty$.


Но тогда существует такая подпоследовательность


$(n_m)_{m=1}^\infty$, что $T_{n_m}\to 0$ п.\,н. при $m\to\infty$.


По лемме 2 среднее


${\bold E}f(\sup\limits_{n\ge 1}\| T_{n} \|)<\infty$,


откуда в силу теоремы Лебега о мажорантной сходимости


${\bold E}f(\sup\limits_{k\ge n_m}\| T_{k} \|)\to 0$


при $m\to\infty$. Итак,


${\bold E}f(\sup\limits_{k\ge n}\| T_{k} \| ) \to 0$ при $n\to \infty $.


\end{pf}





\begin{cornonum} Пусть заданы взвесь ${\bold a}$ и соответствующая


последовательность $\varphi $ и пространство $E$ имеет радемахеровский тип


$p$, $1<p<2$. Если $(X_{k} )$ --- такие независимые симметричные


случайные элементы, что


$\sum\limits^{\infty }_{k=1} {\bold E} \| X_{k} \| ^{p}


/\varphi ^{p} (k)<\infty $, то $T_{n} \to 0$ п.\,н.


\end{cornonum}





Теперь, используя лемму 1 и предложение, осуществим





\begin{pf*}{Доказательство теоремы 1}


Так как $\varphi $ есть правильно меняющаяся


последовательность с показателем $1/p$, то выполняются


условия леммы 1 для всех $r>p$ и $q<p$ (\cite{Sen}, упражнение 2.1 и


теорема 2.8). Обозначим


\begin{gather*}


X'_k= X_{k} I\{\| X_{k} \|


<\varphi (k)\},\quad


X''_k= X_{k} I\{\| X_{k} \| \ge\varphi (k)\},\\


%


T'_n=\sum^{n}_{k=1} a_{k} (n)X'_k,\quad


T''_n=\sum^{n}_{k=1} a_{k} (n)X''_k.


\end{gather*}


Тогда


$X_{k} =X'_k+X''_k$ и


$$


{\bold E}


f(\sup_{n\ge k}\| T_{n} \| )=


{\bold E} f(\sup_{n\ge k}\| T'_n+T''_n\| )


\le С{\bold E} f(\sup_{n\ge k}\| T'_{n}\| )+


С{\bold E} f(\sup_{n\ge k}\| T''_n\| ),


$$


т.\,к. функция $f$ полуаддитивна.





Пространство $Е$ имеет устойчивый тип $p< 2,$ значит, оно


имеет радемахеровский тип $r$ для некоторого $r>p$ (\cite{Pis}, сh.\,4) и


по лемме 1(1)


$$


\sum^{\infty }_{k=1} {\bold E}


\| X'_k\| ^{r} /\varphi ^{r} (k)<\infty.


$$


Следовательно, по предложению $T'_n\to 0$ в $\fs$.





Так как $f\in {\bold G}_{p} $, то существует такое $q<p$,


что $f\in {\bold I}_{q} $. По


лемме 1(2) $\sum\limits^{\infty }_{k=1} {\bold E} \|


X{''}_n\| ^{q} /\varphi ^{q} (k)<\infty $.


Следовательно, по предложению


$T_n\to 0$ в $\fs$.


\end{pf*}





\section{Применения}





Теорема 1 связана с некоторыми задачами


планирования оптимального момента остановки, а именно, с


вопросом об ограниченности


${\bold E} \sup\limits_{n\ge 1}\Bigl(\Bigl|\sum\limits^{n}_{k=1} \xi _{k}


\Bigr|/ \varphi (n)\Bigr)^{r}$ для некоторых


$\varphi (n)$ и $r>0$\linebreak


%//                  ^^^^^^^^


(напр., $\varphi (n)=n^{q}$ \cite{Hof}, \cite{Choi}) и


независимых одинаково распределенных центрированных случайных величин


$\xi _{k} $.


Приведем один результат такого рода для случайных элементов


со значениями в банаховом пространстве.





\begin{thm} Пусть банахово пространство $Е$ имеет устойчивый


тип $p$, $1<p<2$ последовательность


$\varphi =(\varphi (n))$ правильно меняется


с показателем $1/p$, функция $f\in {\bold G}_{p}$, $(X_{k} )$ ---


последовательность независимых симметричных случайных элементов с


одинаково распределенными нормами. Следующие утверждения


эквивалентны{\em :}


\begin{itemize}


\item[(1)] $\sum\limits^{\infty }_{k=1} {\bold P} \{\| X_{1} \|


>\varphi (k)\}< \infty$,


\item[(2)] $S _{n} /\varphi (n)\to 0$ в $\fs$ при $n\to\infty$,


\item[(3)] $ {\bold E} f(\sup\limits_{n\ge 1}\| S _{n} \| /


\varphi (n))< \infty$,


\item[(4)] $ {\bold E} f(\sup\limits_{n\ge 1}\| X_{n} \| /


\varphi (n))< \infty$.


\end{itemize}


\end{thm}





\begin{pf} Импликация (1)$\Longrightarrow$(2) установлена в теореме


1, а (2)$\Longrightarrow$(3) тривиально. Покажем, что


(3)$\Longrightarrow$(4). Так


как функция $f$ полуаддитивна и $\varphi (n)$ возрастает, то


\begin{gather*}


{\bold E} f(\sup_{n\ge 1}


\| X_{n} \| /\varphi (n))= {\bold E} f(\sup_{n\ge 1}


\Bigl\| S_{n} /\varphi (n) - \frac{\varphi(n)}{\varphi (n-1)} S_{n-1} /


\varphi (n)\Bigr\|\le \\


%


\le C{\bold E} f(\sup_{n\ge 1}\| S_{n} \| /\varphi (n))+


C{\bold E} f(\sup_{n\ge 1}\| S_{n-1} \| /


\varphi (n-1))


\le 2C{\bold E} f(\sup_{n\ge 1}\| S_{n} \|/\varphi (n)).


\end{gather*}


Наконец, докажем от противного, что (4)$\Longrightarrow$(1). Предположим,


что 


$$


\sum^{\infty }_{k=1} {\bold P}


\{\| X_{1} \| >\varphi (k)\}= \infty.


$$ 


Тогда для всех $M\in {\bold N} $


$$


\infty =\sum^{\infty }_{k=1} {\bold P}


\{\| X_{1} \| >\varphi (Mk)\}


=\sum^{\infty }_{k=1} {\bold P}


\{\| X_{1} \| >C\varphi(M)\varphi (k)\}


\le \sum^{\infty }_{k=1}


{\bold P} \{\| X_{1} \| /\varphi (k)>C


\varphi(M)\},


$$


где в предпоследнем неравенстве использовалось условие правильной


меняемости $\varphi$. По лемме Бореля-Кантелли


$\sup\limits_{n\ge 1} \| X_{n} \| /\varphi (n)>\varphi (M)$


п.\,н., т.\,е.


$\sup\limits_{n\ge 1}\| X_{n} \| /\varphi (n)=\infty$


п.\,н. Следовательно,


${\bold E}


f(\sup\limits_{n\ge 1}\| X_{n} \| /\varphi (n)) =\infty$,


что противоречит (4).


\end{pf}





В заключение выражаю признательность рецензенту моей


статьи, профессору Т.\,Микошу (Цюрих, Швейцария),


профессору А.\,Адлеру (Чикаго,


США) и профессору Д.\,Шиналю (Люблин, Польша) за полезные


замечания.





\begin{thebibliography}{99} 


 


\bibitem{Kra} 


Красносельский М.А., Рутицкий Я.Б. {\em Выпуклые функции и пространства 


Орлича}. -- М.: ГИФМЛ, 1958. -- 271 с. 


 


\bibitem{Pis} 


Pisier G. {\em Probabilistic methods in the geometry of Banach spaces} // 


Lect. Notes. Math. -- 1986. -- \No\,1206. -- P.\,167--241. 


 


\bibitem{Sz} 


Sztencel R. {\em On boundedness and convergence of some Banach space valued 


random variables} // Probab. and Math. Stat. -- 1981. -- V.\,2. -- 


P.\,83--88. 


 


\bibitem{Azl} 


Азларов Т.А., Володин Н.А. {\em Законы больших чисел для одинаково 


распределенных банаховозначных величин}. -- Теория вероятн. и ее примен. -- 


1981. -- Т.\,XXVI. -- \No\,3. -- С.\,573--580. 


 


\bibitem{Aco} 


Acosta A. de. {\em Inequalities for $B$-valued random vectors and 


applications to the strong law of large numbers} // Ann. Probab. -- 1981. 


-- V.\,9. -- P.\,157--161. 


 


\bibitem{Marc} 


Marcus M., Woyczynsky W. {\em Stable measures and central limit theorem in 


spaces of stable type} // Trans. Amer. Math. Soc. -- 1979. -- V.\,271. -- 


P.\,71--102. 


 


\bibitem{Adl} 


Adler A., Rosalsky A., Taylor R. {\em Strong laws of large numbers for 


weighted sums of random elements in normed linear spaces} // Int. J. Math. 


and Math. Sci. -- 1989. -- V.\,12. -- P.\,509--530. 


 


\bibitem{Mik} 


Mikosh T., Norvaisa R. {\em Strong laws of large numbers for fields of 


Banach space valued random variables} // Probab. Th. Rel. Fields. -- 1987. 


-- V.\,74. -- P.\,241--253. 


 


\bibitem{Tay} 


Taylor R. {\em Stochastic convergence of weighted sums of random elements 


in linear spaces}. -- Lect. Notes. Math. -- 1976. -- \No\,672. -- 


270 p.


 


\bibitem{Vol1} 


Володин А.И. {\em Сходимость взвешенных сумм случайных элементов в 


пространствах типа $p$} // Изв. вузов. Математика. -- 1990. -- 


\No\,10. -- С.\,65-67. 


 


\bibitem{Vol2} 


Volodin A.I. {\em On the law of large numbers for weighted sums of random 


elements in $p$-type spaces} // V Int. Vilnius Conf. on Probab. Th. and 


Math. Stat. Abstracts of Commun. -- 1989. -- V.\,II. -- P.\,216--217. 


 


\bibitem{Sen} 


Сенета Е. {\em Правильно меняющиеся функции}. -- М.: Наука, 1985. -- 142 с. 


 


\bibitem{Sto} 


Stout W. {\em Almost sure convergence}. -- New~York: Acad. Press. -- 1974. 


 


\bibitem{Choi} 


Choi B.D., Sung S.H. {\em On moment conditions for the supremum of normed 


sums}. -- Stochast. Proc. and Appl. -- 1987. -- V.\,26. -- P.\,99--106. 


 


\bibitem{Hof} 


Hoffmann-Jorgensen J., Pisier G. {\em The law of large numbers and the 


central limit theorem in Banach spaces} // Ann. Probab. -- 1976. -- V.\,4. 


-- P.\,587--599. 


 


\bibitem{Gut} 


Gut A. {\em Moments of the maximum of normed partial sums of random 


variables with multidimensional indices} -- Zeitsch. Wahrsch. -- theorie Verw. 


DDR Gebiete. -- 1979. -- Bd.\,46. -- S.\,205--220. 


 


 


\end{thebibliography} 


 


\vskip 2cm 


 


\post{\it Казанский государственный университет}{\it Поступила} 


\post{}{24.02.1995} 


 


\end{document} 


